\documentclass[11pt]{article}
\usepackage[margin=1in]{geometry}
\usepackage{amsmath,amssymb,amsthm}

\newtheorem{theorem}{Theorem}
\newtheorem{lemma}[theorem]{Lemma}
\DeclareMathOperator{\Nm}{N}

\begin{document}

\begin{theorem}\label{thm:classification}
Let $p,q$ be primes. If $p=2^m+q^n$ has two distinct solutions in
positive integers $m,n$, then $(p,q)=(11,3)$.
\end{theorem}

In two distinct solutions the larger exponent of $2$ occurs with the
smaller exponent of $q$. Write
\begin{equation}\label{eq:partitions}
p=2^{c+M}+q^e=2^c+q^{e+N},\qquad c,e,M,N>0.
\end{equation}

\begin{lemma}\label{lem:gaps}
The base $q$ is odd, $M$ is even, and $N$ is odd. There is a positive
integer $A$ such that
\begin{equation}\label{eq:gaps}
Aq^e=2^M-1,\qquad A2^c=q^N-1,\qquad Ap=2^Mq^N-1.
\end{equation}
Moreover, $q^e\le 2^{M/2}+1$.
\end{lemma}

\begin{proof}
Since $p>3$ is prime, it is odd, hence so is $q$. Subtraction gives
$2^c(2^M-1)=q^e(q^N-1)$. Coprimality gives the first two
identities in \eqref{eq:gaps}, and
$Ap+1=(Aq^e+1)(A2^c+1)$ gives the third.

Modulo $3$, the binary summands have the same residue. If $q=3$,
both are congruent to $p$. Otherwise both avoid the nonzero residue
$p$, since the $q$-summands are nonzero. Thus $M$ is even.

If $N$ were even, $X=2^{M/2}q^{N/2}$ would satisfy
$1<X\le p-2$, since both $2^M$ and $q^N$ are at most $p-2$.
But $p$ divides $(X-1)(X+1)=Ap$, whose factors are strictly between
$0$ and $p$.

The factors $2^{M/2}-1$ and $2^{M/2}+1$ are coprime, so one
contains the full power $q^e$.
\end{proof}

\begin{lemma}\label{lem:three}
If $q=3$, then $p=11$.
\end{lemma}

\begin{proof}
Since $N$ is odd, $A2^c=3^N-1\equiv2\pmod4$, so $c=1$.
If $e\ge2$, then $2^M\equiv1\pmod9$. The order of $2$ modulo $9$
is $6$. Since $2^6\equiv1\pmod7$, we get $7\mid2^M-1$ and $7\mid A$. Thus
$3^N\equiv1\pmod7$, contrary to $N$ odd because the order of $3$
modulo $7$ is $6$.

Thus $e=1$, and $(3^{(N+1)/2}-1)(3^{(N+1)/2}+1)=2^{1+M}$.
The factors are powers of $2$ differing by $2$, hence $2$ and $4$.
Therefore $p=2+3^2=11$.
\end{proof}

\begin{lemma}\label{lem:exponents}
If $q>3$, then $c$ and $e$ are even and $2^c\le q-1$.
\end{lemma}

\begin{proof}
Since $M$ is even, $3\mid 2^M-1$, hence $3\mid A$.
It follows that $q^N\equiv1\pmod3$, so $q\equiv1\pmod3$ by oddness
of $N$. Now $p\equiv 2^c+1\pmod3$ and $3\nmid p$, so $c$ is even.

The quotient $(q^N-1)/(q-1)$ is odd, as is $A$.
Thus $c=v_2(q-1)$, giving $2^c\le q-1$.
If $e$ were odd, $e+N$ would be even. Set
$u=2^{(c+M)/2}$ and $v=q^{(e+N)/2}$. By \eqref{eq:partitions},
$v^2-u^2=q^e-2^c>0$. Lemma~\ref{lem:gaps} then gives
\[
2u+1\le v^2-u^2<q^e\le 2^{M/2}+1,
\]
contrary to $u\ge 2^{M/2}$.
\end{proof}

\begin{lemma}\label{lem:norm}
Let $D>1$ be an odd integer and $q$ an odd prime. If $x,y$ are
even integers, $q\nmid x$, and
\[
x^2+D=q^n,\qquad y^2+D=q^e
\]
for positive integers $n,e$, then $n\equiv e\pmod2$.
\end{lemma}

\begin{proof}
Choose the sign of $y$ so that $q\nmid x+y$, since $q\nmid2x$.
In the ring $R=\mathbb Z[w]$, where $w=\sqrt{-D}$, put
$\alpha=x+w$, $\beta=y+w$, and $H=q^{ne}$. Then
\[
\alpha^e\bar\alpha^e=\beta^n\bar\beta^n=H.
\]
In the quotient ring $R/(H)$, the sum $\bar\alpha+\beta=x+y$ is a unit. Hence
$\bar\alpha^e$ and $\beta^n$ generate the unit ideal there.
Both annihilate $\alpha^e\bar\beta^n$, so
$\alpha^e\bar\beta^n=Hz$ for some $z\in R$.
Taking norms gives $\Nm(z)=1$, so $z=\pm1$ because $D>1$.

In the quotient ring $R/2R$, both $\alpha$ and $\bar\beta$ reduce to $w$, and
$Hz$ reduces to $1$. Thus $w^{e+n}=1$. Since $w^2=1$ and
$w\ne1$ in $R/2R$, $e+n$ is even.
\end{proof}

\begin{proof}[Proof of Theorem~\ref{thm:classification}]
By Lemmas~\ref{lem:gaps} and~\ref{lem:three}, it suffices to exclude
$q>3$. Set $x=2^{(c+M)/2}$, $y=2^{c/2}$, and $D=q^e-2^c$.
By Lemma~\ref{lem:exponents}, $x,y$ are even and
$D\ge q^2-q+1>1$. The integer $D$ is odd, and
\eqref{eq:partitions} gives $x^2+D=q^{e+N}$ and $y^2+D=q^e$.
Since $q\nmid x$, Lemma~\ref{lem:norm} contradicts the oddness of $N$.
The two solutions at $(11,3)$ are $11=2^3+3=2+3^2$.
\end{proof}

\end{document}
